\documentclass[11pt]{article}
\usepackage[margin=1in]{geometry}
\usepackage{amsmath,amssymb,amsthm}
\usepackage{hyperref}

\newtheorem{definition}{Definition}
\newtheorem{proposition}{Proposition}

\title{\textbf{CRC-S (Abelian): Cycle-Residual Consistency on Cellular Sheaves with Abelian Connection}}
\author{Beexly Research}
\date{November 2026}

\begin{document}
\maketitle

\begin{abstract}
We formalize Cycle-Residual Consistency (CRC-S) for finite connected graphs equipped with an abelian structure group $A \in \{U(1), \mathrm{SO}(2)\}$. We relate cycle residuals to discrete holonomy, characterize the global section space $\dim H^0$, distinguish minimal angle-edit surgery from undirected minimum feedback edge sets, and define Relative CRC against a system-of-record reference to separate self-consistent fabrication from statistical noise.
\end{abstract}

\section{Setup and Definitions}
Let $G = (V, E)$ be a finite connected graph, $A$ a compact abelian structure group (here $A = U(1)$ or $\mathrm{SO}(2)$), and $\{\theta_e\}_{e \in E}$ edge angles defining a discrete connection. Let $\{\gamma_i\}$ be an oriented cycle basis of $G$ with orientation incidence $\varepsilon_e \in \{-1, 0, +1\}$.

Define the \emph{cycle residual} on cycle $\gamma_i$ by
\begin{equation}
  r_i = \left( \sum_{e \in \gamma_i} \varepsilon_e \theta_e \right) \bmod 2\pi.
\end{equation}

\begin{definition}[CRC-flatness]
The connection $\{\theta_e\}_{e \in E}$ is \textbf{CRC-flat} if $r_i = 0$ for all cycles $\gamma_i$ in the basis.
\end{definition}

\section{Main Propositions}

\begin{proposition}[Holonomy and Cycle Residuals]
For an abelian structure group $A$, the holonomy $\mathrm{Hol}_{\gamma_i} = I$ if and only if $r_i = 0$. Consequently, a non-trivial holonomy flag occurs if and only if at least one cycle residual $r_i$ is non-zero.
\end{proposition}

\begin{proposition}[Global Sections]
Suppose the connection is CRC-flat and the stalks are $\mathbb{R}^d$ with compatible restrictions after gauge fixing $\theta = 0$. Then the zeroth cohomology dimension satisfies
\begin{equation}
  \dim H^0 = d \cdot \#\{\text{connected components of } G\}.
\end{equation}
Conversely, if $r_i \neq 0$ for some cycle $\gamma_i$ in the pure frame case, then $\dim H^0 = 0$ (no non-trivial global parallel frame exists).
\end{proposition}

\begin{proposition}[Minimal Surgery vs.\ Feedback Edge Set]
Editing a subset $S \subset E$ of edge angles to new values restores CRC-flatness if and only if the induced cycle residuals vanish. A minimal-cardinality edit set $S$ corresponds to a minimum-weight edit solution for the linear system
\begin{equation}
  C\theta \equiv 0 \pmod{2\pi},
\end{equation}
where $C$ is the cycle-edge incidence matrix. In general, $|S|$ is \textbf{not} equal to the size of a minimum feedback edge set (FES) of the underlying undirected graph $G$.
\end{proposition}

\begin{proposition}[Relative CRC]
Let $x$ denote stalk values across peripheral systems excluding a System of Record (SoT), with CRC-flat maps among them and sheaf Dirichlet energy $E(x) = 0$. Let $x_{\mathrm{ref}}$ denote SoT reference values. When $\|x - x_{\mathrm{ref}}\|$ is large, the configuration is classified as \textbf{self\_consistent\_wrong\_vs\_ledger}. Value-difference checks alone, without the flatness check, cannot separate such coordinated fabrication from unstructured statistical noise.
\end{proposition}

\section{Empirical Fixtures and Scope}
\begin{itemize}
  \item \textbf{Penrose triangle fixture} ($r = \pm \pi/2$): triggers structural holonomy repair class.
  \item \textbf{Schema-real KYC CSV fixture}: separates relative fraud (\texttt{E001}), aligned (\texttt{E002}), and statistical noise (\texttt{E003}).
\end{itemize}

\paragraph{Non-claims.} This note does not claim non-abelian structure groups, continuum Yang--Mills limits, or production fraud-rate reductions.

\section{Citation Spine}
Ambrose--Singer (holonomy generated by curvature); Hansen--Ghrist (cellular sheaf Laplacians); Wilson (plaquette as discrete curvature); classical discrete exterior calculus ($d^2 = 0$).

\end{document}
